\section{The local prolate limit over the rational field}

This section proves the fixed-frequency comparison which connects the
Hardy model spaces to the finite arithmetic prolate images.
The analytic estimates needed from the cutoff are precisely the
concentration bound, inverse Gram estimate, and theta identities
listed before Theorem~\ref{thm:allzero-lower}, together with the
elementary critical-line bound
\begin{equation}\label{eq:local-theta-graph}
 \|E_\chi f\|_{L^2(dt)}\leq C_\chi\|f\|_{H_D^{5/16}}.
\end{equation}
Here $H_D^s$ denotes ordinary $H^s(\mathbb R)$ after the unitary
$(Uf)(t)=e^{t/2}f(e^t)$.  The closed theta map is its Mellin multiplier
$L(1/2+i\tau,\chi)$, corrected by the principal polar subtraction.
All these inputs are independent of the location of nontrivial zeros.
Write $(-1)^\epsilon=\chi(-1)$ and let $\mathcal F_\epsilon$
be the raw Fourier transform restricted to this parity space.
When importing the abstract endpoint notation, put
$Q_L=Q_{\chi,e^L}$.

\begin{lemma}\label{lem:local-prolate-interior}
Let $p\in C_c^\infty((-1,1))$ be even, equal to one near zero, and
$0\leq p\leq1$. Set
\[
 P^p=p(u/\Lambda),\qquad
 B^p=\mathcal F_\epsilon^{-1}p(\xi/R)\mathcal F_\epsilon.
\]
For every $A>0$ there is a constant, depending only on
$\chi,p,A$, such that
\[
 \|E_\chi(1-\Pi_\Lambda)P^pB^pP^p\|_{L^2(du)\to L^2(dt)}
                                    \leq C_{\chi,p,A}c^{-A}.
\]
In particular, if $T_L$ is a raw Hilbert--Schmidt column operator with
$\|T_L\|_2=O((1+L)^d)$, then for every fixed $A>0$,
\[
 \|E_\chi(1-\Pi_\Lambda)P^pB^pP^pT_L\|_2=O(c^{-A}).
\]
The same graph approximation holds for each fixed finite family of
smooth profiles supported strictly inside the unit raw time interval,
dilated to radius $\Lambda$, and for its Fourier-dual family dilated
to frequency radius $R$.
\end{lemma}

\begin{proof}
Choose $b<1$ containing the support of $p$.
Abbreviate $P=P_\Lambda$ and $B=B_R$.
The Fourier kernel of $P^p$ is
$\Lambda\widehat p(\Lambda(\xi-\eta))$.
Between $|\eta|\leq bR$ and $|\xi|>R$ its arguments are separated
by at least $(1-b)c$.  Schur's test and the Schwartz decay of
$\widehat p$ give a band-tail operator norm $O_N(c^{-N})$ for every N.
The time tail of $P^pB^pP^p$ is zero.
The spectral theorem for the prolate concentration operator, using a
fixed upper bound below one for the first omitted eigenvalue, gives
\[
 \|(1-\Pi_\Lambda)f\|^2
   \leq C\bigl(\|(1-P)f\|^2+\|(1-B)f\|^2\bigr).
\]
Thus the raw projection error has every polynomial rate, also after
summing squared column norms. All constants at this stage belong
to the raw operator estimate and are independent of the number
or choice of columns.

For $0<s<1/2$, the selected prolate vectors satisfy
$\|\Pi_\Lambda f\|_{H_D^s}\leq C_s(1+c)^s\|f\|$.
One proof differentiates the band extension in the time interval,
whose logarithmic derivative costs at most $C(1+c)$, then uses the
fractional Sobolev estimate for extension by zero across a single
endpoint.  That estimate follows by integrating by parts in the
Fourier transform and integrating its boundary term
$|\tau|^{2s-2}$, which is integrable exactly for $s<1/2$.
The smooth-cutoff vector has the same coarse bound, since its raw
band derivative costs at most $2\pi R$ and $u\leq\Lambda$ on its
time support.  Interpolating raw smallness with the $H_D^{3/8}$
bound gives smallness in $H_D^{5/16}$, and
\eqref{eq:local-theta-graph} proves the assertion.
For fixed profiles the complementary raw tail is Schwartz at the
scale c, and the same reasoning applies.  Raw Fourier transformation
reverses the logarithmic dilation generator and preserves its
Sobolev norms, giving the dual assertion.
\end{proof}

Here $H^2$ is the Hardy space on $\Re s>1/2$, with norm
\[
 \|F\|_{H^2}^2=\sup_{\sigma>1/2}
       \int_\R|F(\sigma+i\tau)|^2\frac{d\tau}{2\pi}<\infty.
\]
Boundary values on $s=1/2+i\tau$ identify it with the Fourier
transforms of $\mathcal H_-$. The inverse Fourier notation below
uses these boundary values.

\begin{lemma}\label{lem:local-halfline-closure}
In a nonprincipal sector the closure of theta images of smooth raw
profiles supported in $0<u<1$ is
$M_+=\mathcal F^{-1}(B_\chi H^2)$ inside $\mathcal H_-$, where
$B_\chi$ is the Blaschke product of the zeros of $L(s,\chi)$ in
$\Re s>1/2$.  The same assertion holds in the principal sector
when both raw mean and value at zero are required to vanish.
The Fourier-dual family has the reflected closure $M_-$ at the
opposite endpoint. Its complementary model space records the
left-half-strip zeros of the original character.
\end{lemma}

\begin{proof}
The nonprincipal statement is Burnol's closure theorem; its
number-field proof uses the fixed standard finite local vectors and
arbitrary smooth archimedean profiles supported in $(0,1)$.
It can also be seen directly by the invariant-subspace argument
which clarifies the principal condition.
The closed image is invariant under contractions of raw scale and
is supported in the target half-line.  The Beurling theorem identifies
it with an inner multiple of the half-line Hardy space.
For the nonprincipal family its Mellin transforms are
$L(s,\chi)\mathcal M\phi(s)$.
For the principal family take
$\psi=(u\partial_u+1)\phi$, where $\phi$ is smooth and compactly
supported in $(0,1)$.  Then $I\psi=\psi(0)=0$, and the transform is
$(1-s)\zeta(s)\mathcal M\phi(s)$.
At any prescribed point one may choose a sufficiently narrow bump
$\phi$ with nonzero Mellin transform.  Hence the common zeros in
the Hardy half-plane are exactly the indicated nontrivial zeros;
the principal factor is regular and nonzero at $s=1$ after pole
cancellation.  Each generator continues analytically across the
finite boundary line, so a common singular inner measure there is
impossible.  A common exponential inner factor would force all
target functions to vanish on a fixed interval immediately below
zero in logarithmic coordinates.  This is impossible: choose raw
profiles supported arbitrarily close to $u=1$; for $x>1/2$ their
theta sum has only its $n=1$ term.  In the principal case the
differentiated bump still has mean zero and is nonzero arbitrarily
close to that endpoint.  There is therefore no common delay.
This identifies the inner factor exactly.
The primitive theta--Fourier identity, with $R=\Lambda/q$, puts the
dual endpoint at $1/(qR)=\Lambda^{-1}$ and exchanges the character
with its conjugate.  The functional equation identifies its
right-hand zeros with the original left-hand zeros, with their
multiplicities.  Reflection proves the last assertion.
\end{proof}

\begin{theorem}\label{thm:rational-local-limit}
For each fixed
primitive Dirichlet character, including the constrained principal
sector, and every $a\in C_c^\infty(\mathbb R)$,
\[
 \lim_{\Lambda\to\infty}\operatorname{Tr}(D_LC_aC_a^*)
                   =\int |a(\tau)|^2\,d\mathfrak b_\chi(\tau).
\]
No hypothesis on the location or multiplicity of the zeros is used.
The positive defect measures consequently converge vaguely to
$\mathfrak b_\chi$.  This theorem by itself asserts no convergence
against noncompact frequency tests.
\end{theorem}

\begin{proof}
Choose a fixed compact interval whose interior contains
$\operatorname{supp}a$ and whose endpoints are not ordinates of
critical-line zeros.  Let $P_{0,L}$ project onto the finite span of
the window jets of the critical-line zeros in this interval, and set
$W_L=S_L-P_{0,L}$.  A pole outside the support of a has zero coefficient
and causes no difficulty.  Define the compactly supported distribution
\[
 b(\tau)=\frac{a(\tau)}{L(1/2+i\tau,\chi)}
\]
by prescribing $(\tau-\gamma-i0)^{-j}$ at each real pole.
Its inverse Fourier transform $\kappa$ is smooth.  At $+\infty$ it
and all its derivatives decrease faster than every power; at
$-\infty$ it is a finite exponential polynomial, of degree below
the corresponding zero multiplicities, plus a Schwartz remainder.
Indeed subtract each principal part with a compact smooth cutoff.
The remaining distribution is a smooth compactly supported function.
The inverse transform of a prescribed pole is a constant times
$1_{t<0}t^{j-1}e^{-i\gamma t}$; convolving with the Schwartz inverse
transform of its cutoff gives precisely those two tail descriptions.
Off-line zeros do not create real distributional poles and require
no modification of this argument.

Put $T_L=U^{-1}C_bW_L$.  The zero-jet conditions cancel every
exponential-polynomial term in the output tail below $-L$.
The other tail is Schwartz.  On the middle interval the kernel and
each fixed derivative have polynomial growth in L.
Integrating the squared kernel on these three regions proves
\[
 \|UT_L\|_{\mathcal S_2(\mathcal H,H^j)}=O_j((1+L)^{d_j})
 \quad(j\geq0),\qquad E_\chi T_L=C_aW_L.
\]
The last identity follows by cancellation of each distributional
pole in multiplication by the analytic L-factor and the closed
multiplier realization.

The ordinary raw Fourier transform has an equally useful expression.
Write $\epsilon\in\{0,1\}$ for the parity of $\chi$,
$\tau(\chi)=\sum_{a\bmod q}\chi(a)e^{2\pi ia/q}$, and
$(J_qf)(t)=f(-t-\log q)$. With the Fourier convention fixed above,
Poisson summation gives the exact identities
\[
 E_\chi=\frac{\tau(\chi)}{\sqrt q}J_qE_{\bar\chi}\mathcal F,
 \qquad
 E_\chi\mathcal F
   =\frac{(-1)^\epsilon\tau(\chi)}{\sqrt q}J_qE_{\bar\chi}.
\]
The second follows from the first and
$\mathcal F^2=(-1)^\epsilon$ on the parity space. Thus the
reflected inverse multiplier, with
$\kappa_{\bar\chi}=(-1)^\epsilon\tau(\bar\chi)/\sqrt q$, is
\[
 b^\#(\tau)=\kappa_{\bar\chi}q^{-i\tau}
                 \frac{a(-\tau)}{L(1/2+i\tau,\bar\chi)}.
\]
This is compactly supported as a distribution. Changing its real
pole prescriptions back to the negative-time prescriptions adds
only delta derivatives, which vanish on the reflected $W_L$.
Consequently $\mathcal FT_L$ has the same kernel description,
polynomial bounds, and one-sided Schwartz tail, with the fixed
translation $\log q$. No condition on off-line zeros enters here.

Use $P^p,B^p$ from Lemma~\ref{lem:local-prolate-interior}.
After moving the raw upper endpoint to zero and the input endpoint
to zero, $(1-P^p)T_L$ has, before the finite jet correction, kernel
\[
       (1-p(e^t))\kappa(t-s)1_{[-2L,0]}(s).
\]
Its output coordinate is bounded below by a fixed number, and its
input coordinate is nonpositive.  The positive Schwartz tail of
$\kappa$ therefore makes this kernel converge in Hilbert--Schmidt
norm, with any fixed output logarithmic derivative, to the same
kernel with $1_{(-\infty,0]}(s)$.
Ordinary raw derivatives insert factors $e^{-t}$, bounded on this
output support, so the same conclusion holds for them.
The finite jet correction tends to zero in these norms: each
normalized translated jet has size $O(L^{-1/2})$ times a fixed-degree
polynomial in $s/L$, and the quadrant kernel has all integrable
moments in s.  The Fourier-dual upper tail has the identical
convergence, with only the fixed translation $\log q$ in its kernel.

These derivative bounds also imply
\[
 \|(1-B^p)(1-P^p)T_L\|_{2,H_D^j}
 +\|(1-P^p)(1-B^p)T_L\|_{2,H_D^j}=O_{j,N}(c^{-N})
\]
for arbitrary fixed $j,N$.  For the first term, the rescaled
time-boundary kernel has every ordinary raw derivative bounded in
Hilbert--Schmidt norm, so its Fourier tail above frequency c has
every polynomial decay rate.  The dual argument gives the second.
Derivatives of the smooth logarithmic cutoffs remain uniformly
bounded, so the estimate also holds in the displayed graph norms.

In a nonprincipal sector, raw time support and the primitive Poisson
identity imply the exact decomposition
\[
 \begin{split}
 S_LE_\chi T_L-E_\chi P^pB^pP^pT_L
 ={}&S_LE_\chi(1-P^p)T_L
     +S_LE_\chi(1-B^p)P^pT_L\\
    &+1_{t\leq L}E_\chi(1-P^p)B^pP^pT_L.
 \end{split}
\]
For verification, expand the first raw time cutoff and then the
band cutoff.  The band image has no target component below $-L$,
whereas the final time-supported image has none above L; these
two support statements give exactly the last term.
The cross estimates permit replacing $P^pT_L$ by $T_L$ in the
second term and $B^pP^pT_L$ by $P^pT_L$ in the last, with
$o_{\mathcal S_2}(1)$ error after applying the theta map.
The first and last terms then have fixed translated upper-endpoint
Hilbert--Schmidt limits; the second has a lower-endpoint limit by
the theta--Fourier identity.  The far opposite target cutoff tends
strongly to the identity on each fixed Hilbert--Schmidt limit.

There is also the input contribution $S_LC_a(1-S_L)$.
Its kernel is Schwartz.  Splitting the input above L and below $-L$
gives two quadrant kernels, each converging after translation to a
fixed Hilbert--Schmidt operator on the appropriate half-line.
Combining these with the previous boundaries yields
\eqref{eq:endpoint-decomposition}, with input projection $P_{0,L}$
and bulk $F_L=E_\chi P^pB^pP^pT_L$.
The bulk has polynomial Hilbert--Schmidt norm because smooth raw
time and frequency multipliers are uniformly bounded on every
fixed logarithmic Sobolev space.  Its distance to the finite theta
image tends to zero by Lemma~\ref{lem:local-prolate-interior}.

The principal sector requires an explicit correction, which we now
describe.  Choose a smooth even $\beta$ supported in
$1/4<|u|<1/2$, with $I\beta=1$.  Put
\[
 b_{\Lambda,+}=D_\Lambda\beta,\qquad
 b_{\Lambda,-}=\tfrac12\mathcal F D_\Lambda\beta,
 \qquad (D_r f)(u)=r^{-1/2}f(u/r).
\]
For the normalized functionals $I/\sqrt\Lambda$ and
$\operatorname{ev}_0/\sqrt\Lambda$, this pair has identity
constraint matrix.  If $A_L=P^pB^pP^pT_L$, let
$r_+=I(A_L)/\sqrt\Lambda$, $r_-=A_L(0)/\sqrt\Lambda$, and
$A_L^0=A_L-b_{\Lambda,+}r_+-b_{\Lambda,-}r_-$.
Its two constraints vanish.
The upper polar tail of a raw time-supported vector is
$-I\psi\,e^{-t/2}$, and the lower tail of a band vector is
$-\psi(0)e^{t/2}/2$. To see the cancellation explicitly, write
$A=P^p$, $B=B^p$, $g=ABA T_L$, and $P_+=1_{t\leq L}$.
Let $R_0$ denote the three-term residual displayed above. Expansion gives
\[
 R_0=S_LE_1T_L+(P_+-S_L)E_1BAT_L-P_+E_1g,
\]
and hence
\[
 S_LE_1T_L-E_1g
 =R_0+Ig\,e^{-t/2}1_{t>L}
       +\tfrac12g(0)e^{t/2}1_{t<-L}.
\]
Here $(BAT_L)(0)=g(0)$ because the time multiplier equals one
near zero. The constraints of $b_{\Lambda,+}$ and
$b_{\Lambda,-}$ show that their upper and lower tails,
respectively, cancel precisely these two terms. Thus
$S_LE_1T_L-E_1A_L^0$ is $R_0$ plus
\[
 1_{t\leq L}E_1 b_{\Lambda,+}r_+
       +1_{t\geq-L}E_1 b_{\Lambda,-}r_-.
\]

Both rows have fixed translated $L^2$ limits.
For $r_+$, moving $B^p-1$ to the normalized smooth integration
test shows that it differs by every polynomially small raw error
from $\Lambda^{-1/2}I((P^p)^2T_L)$.
After scaling, this is integration of $\kappa(t-s)$ against
$e^{t/2}p(e^t)^2$, for $s\leq0$.
The weight has an exponential negative tail and bounded upper
support; the same quadrant argument proves the unprojected row's
weighted bounds and the projected row's unweighted $L^2$ convergence.
No uniform polynomial moment is asserted after input projection.
For $r_-$, evaluate the band convolution at
zero and use the dual kernel, with weight $e^{t/2}p(e^t)$.
The two output polar profiles are fixed after translation and lie
in $L^2$ on their indicated half-lines.  Hence the extra terms are
fixed Hilbert--Schmidt boundaries of exactly the required type.
The corrected bulk is approximated by the constrained selected
space as follows. The projection $\Pi_\Lambda$ here is the
unconstrained projection. Project its two correction profiles; their
normalized constraint matrix tends to the identity by the graph
approximation, the compact-support mean estimate, and the local
Sobolev estimate for value at zero, justified as follows. For the
projection error $F$, raw $L^2$ smallness has every polynomial
rate, while the ordinary derivative on $(-1,1)$ has polynomial
growth by the band extension. The inequality
\[
 |f(0)|^2\leq C\|f\|_{L^2(-1,1)}
 (\|f'\|_{L^2(-1,1)}+\|f\|_{L^2(-1,1)})
\]
and Cauchy--Schwarz over columns give the required smallness
of the evaluation row. No evaluation claim is made for a
fractional dilation Sobolev space alone. Subtract the two projected
profiles multiplied by the inverse constraint matrix.  This restores
both constraints exactly at a theta graph cost tending to zero.
Thus the same endpoint decomposition holds in the principal sector.

It remains to check the abstract endpoint hypotheses.
The critical jet Gram computation gives
$\|(1-Q_L)S_LC_aP_{0,L}\|_2^2\to\sigma_a$ and the translated
input projections tend strongly to zero.
Lemma~\ref{lem:local-halfline-closure} and the fixed-profile graph
approximation show that every vector in $M_\pm$ is approximated at
its endpoint by the finite theta images.
The finite off-line packet estimates of
Theorem~\ref{thm:allzero-lower} show that the complementary model
vectors are asymptotically annihilated by $Q_L$.
Those packet spans are dense in the Blaschke model spaces: finite
Blaschke products give their reproducing kernels with derivatives,
and the increasing products exhaust the model space.
The model traces of $C_aC_a^*$ are finite by the Poisson-kernel
summability already proved.  The endpoint criterion therefore gives
\[
 \sigma_a+\|P_{K_+}C_a\|_2^2+\|P_{K_-}C_a\|_2^2
                 =\int|a|^2\,d\mathfrak b_\chi.
\]
This proves the asserted local limit.  Polarization gives convergence
against smooth compact frequency tests.  A smooth squared test equal
to one on any given compact set bounds its masses uniformly; uniform
approximation by smooth tests then gives vague convergence.
\end{proof}

\paragraph{Limit of this argument.}
All frequency intervals and their real pole sets were fixed before
letting $\Lambda$ grow.  The constants in the smooth part of
$a/L$ may be large when an off-line zero approaches the frequency
axis.  No bound uniform in that distance was proved above.  A
separate uniform tail argument is therefore required to pass to
the Fourier transforms of arbitrary compact norm tests.
Theorem~\ref{thm:unconditional-frequency-tightness} proves that
moving-window estimate in the next section.
